problem:  
Let \((M^{2n+1}, \mathcal{D}, J)\) be a compact strictly pseudoconvex CR manifold underlying a Sasakian structure. Consider the **Sasaki cone** \(\mathfrak{t}^+\) of Reeb vector fields \(\xi\) compatible with \((\mathcal{D}, J)\), and for each \(\xi \in \mathfrak{t}^+\) let \((\eta_\xi, \xi, \Phi_\xi, g_\xi)\) denote the corresponding Sasakian structure, normalized so that the total contact volume \(\int_M \eta_\xi \wedge (d\eta_\xi)^n = 1\).  

For each such \(g_\xi\), let \(h_{\mathrm{SR}}(\xi)\) denote the topological entropy of the **horizontal geodesic flow** determined by the sub-Riemannian structure \((\mathcal{D}, g_\xi|_{\mathcal{D}})\) on the unit horizontal cotangent bundle \(S^*\mathcal{D}\).

**Question:**  
Does the function \(\xi \mapsto h_{\mathrm{SR}}(\xi)\) on the Sasaki cone attain its global minimum precisely at the quasi‑regular Reeb vector fields (those generating circle actions), and are all minimizing directions necessarily quasi‑regular?  
Equivalently, can an irregular Reeb field yield dynamically simplest (i.e. entropy-minimizing) horizontal behavior, or does entropy rigidity enforce quasi‑regularity of the Reeb foliation?

Why is it a "good" problem:  
This problem asks whether a **dynamical extremal principle**—minimization of the horizontal geodesic entropy—detects the arithmetic regularity of Sasakian Reeb fields. It is mathematically sound: the Sasaki cone parametrizes all compatible Reeb vector fields with fixed CR structure; normalization fixes volume, ensuring \(h_{\mathrm{SR}}(\xi)\) is well defined up to scale.  

Conceptually, it forges a new link among three major frameworks:
1. **Sasakian geometry and its moduli space** (variation of \(\xi\) inside the Sasaki cone),  
2. **Sub‑Riemannian geometry and horizontal dynamics** (entropy of geodesic flow restricted to the contact distribution), and  
3. **Ergodic and metric rigidity theory** (entropy minimization and symmetry phenomena known in Riemannian settings).

Entropy rigidity results are well studied for negatively curved Riemannian manifolds and locally symmetric spaces, where lowest entropy characterizes maximal symmetry. This problem asks whether a parallel principle survives in the sub‑Riemannian Sasakian world, where quasi‑regularity represents “maximal transverse symmetry.”  

The statement itself is simple—“is the dynamically simplest Sasakian flow necessarily the geometrically most regular?”—yet profoundly deep and presently unresolved. Any progress would open a new pathway between contact/CR geometry, sub‑Riemannian analysis, and dynamical entropy theory, thus qualifying as a truly “good” research‑level mathematical problem.